\begin{thebibliography}{00}
\bibitem{c1} H. Zhou, S. Zhang, J. Peng, S. Zhang, J. Li, H. Xiong, and W. Zhang, ``Informer: Beyond efficient transformer for long sequence time-series forecasting,'' in \textit{Proceedings of the AAAI Conference on Artificial Intelligence}, 2021, pp. 11106--11115.
\bibitem{c2} H. Wu, J. Xu, J. Wang, and M. Long, ``Autoformer: Decomposition transformers with auto-correlation for long-term series forecasting,'' in \textit{Advances in Neural Information Processing Systems}, 2021, pp. 22419--22430.
\bibitem{c3} Y. Nie, N. H. Nguyen, P. Sinthong, and J. Kalagnanam, ``A time series is worth 64 words: Long-term forecasting with transformers,'' in \textit{International Conference on Learning Representations}, 2023.
\bibitem{c4} J. Wu, H. Wu, and M. Long, ``TimesNet: Temporal 2D-variation modeling for general time series analysis,'' in \textit{International Conference on Learning Representations}, 2023.
\bibitem{c5} B. Lim, S. O. Arik, N. Loeff, and T. Pfister, ``Temporal fusion transformers for interpretable multi-horizon time series forecasting,'' \textit{International Journal of Forecasting}, vol. 37, no. 4, pp. 1748--1764, 2021.
\bibitem{c6} S. Hochreiter and J. Schmidhuber, ``Long short-term memory,'' \textit{Neural Computation}, vol. 9, no. 8, pp. 1735--1780, 1997.
\bibitem{c7} K. Cho, B. van Merrienboer, C. Gulcehre, D. Bahdanau, F. Bougares, H. Schwenk, and Y. Bengio, ``Learning phrase representations using RNN encoder-decoder for statistical machine translation,'' in \textit{Proceedings of the 2014 Conference on Empirical Methods in Natural Language Processing}, 2014, pp. 1724--1734.
\bibitem{c8} T. Chen and C. Guestrin, ``XGBoost: A scalable tree boosting system,'' in \textit{Proceedings of the 22nd ACM SIGKDD International Conference on Knowledge Discovery and Data Mining}, 2016, pp. 785--794.
\bibitem{c9} L. Breiman, ``Random forests,'' \textit{Machine Learning}, vol. 45, no. 1, pp. 5--32, 2001.
\bibitem{c10} S. Xingjian, Z. Chen, H. Wang, D. Y. Yeung, W. K. Wong, and W. c. Woo, ``Convolutional LSTM network: A machine learning approach for precipitation nowcasting,'' in \textit{Advances in Neural Information Processing Systems}, 2015, pp. 802--810.
\bibitem{c11} M. R. I. Sheikh, M. S. Islam, and M. S. Rahman, ``Power sector of Bangladesh: Present status and future challenges,'' in \textit{International Conference on Electrical, Computer and Communication Engineering}, 2019, pp. 1--6.
\bibitem{c12} A. S. M. M. H. Rashid, M. A. A. Mahmud, and M. M. Rahman, ``An automated load shedding management system using smart meters in smart grid,'' in \textit{International Conference on Informatics, Electronics \& Vision}, 2016, pp. 53--58.
\bibitem{c13} Y. Bengio, P. Simard, and P. Frasconi, ``Learning long-term dependencies with gradient descent is difficult,'' \textit{IEEE Transactions on Neural Networks}, vol. 5, no. 2, pp. 157--166, 1994.
\bibitem{c14} R. J. Hyndman and G. Athanasopoulos, \textit{Forecasting: principles and practice}. OTexts, 2018.
\bibitem{c15} R. S. Michalski, ``A theory and methodology of inductive learning,'' \textit{Machine Learning: An Artificial Intelligence Approach}, vol. 1, pp. 83--134, 1983.
\bibitem{c16} Open-Meteo, ``Open-meteo weather api,'' https://open-meteo.com, 2024.
\bibitem{c17} J. Chung, C. Gulcehre, K. Cho, and Y. Bengio, ``Empirical evaluation of gated recurrent neural networks on sequence modeling,'' in \textit{NIPS Deep Learning Workshop}, 2014.
\bibitem{c18} M. S. H. Lipu, M. A. Hannan, T. F. Karim, and A. Hussain, ``Intelligent grid monitoring and load forecasting for Bangladesh power sector using machine learning,'' \textit{Energy Reports}, vol. 9, pp. 240--252, 2023.
\bibitem{c19} A. Adadi and M. Berrada, ``Peeking inside the black-box: A survey on Explainable Artificial Intelligence (XAI),'' \textit{IEEE Access}, vol. 6, pp. 52138--52160, 2018.
\bibitem{c20} G. Lai, W.-C. Chang, Y. Yang, and H. Liu, ``Modeling long- and short-term temporal patterns with deep neural networks,'' in \textit{Proceedings of the 41st International ACM SIGIR Conference on Research and Development in Information Retrieval}, 2018, pp. 95--104.
\end{thebibliography}
